The prototype contains six main stages: trace processing, approximate map matching, feature extraction, profile construction, candidate generation, and route ranking. Figure~\ref{fig:pipeline} summarizes the data flow.

\begin{figure*}[t]
  \centering
  \fbox{\begin{minipage}{0.94\textwidth}
  \centering
  OSM public GPS signals / Porto taxi trajectories
  $\rightarrow$ temporal-spatial segmentation
  $\rightarrow$ approximate map matching
  $\rightarrow$ route feature extraction
  $\rightarrow$ dynamic profile construction
  $\rightarrow$ OSM candidate generation
  $\rightarrow$ ranking and evaluation
  \end{minipage}}
  \caption{End-to-end framework. Public movement signals are converted into reconstructed route references and pseudo-history profiles. For a new origin-destination query, the system generates OSM candidates and ranks them with baseline, profile, prompt, or hybrid methods.}
  \Description{A pipeline from public movement signals through segmentation, map matching, feature extraction, profile construction, candidate generation, ranking, and evaluation.}
  \label{fig:pipeline}
\end{figure*}

\subsection{OSM Trace Processing}
The pipeline probes OSM public GPS trackpoints in a Toronto bounding box, parses GPX files, and builds temporal/spatial pseudo-segments. Segments split when timestamp gaps or spatial jumps exceed configured thresholds. Useful segments are retained when they have sufficient distance and plausible speed. These segments are explicitly treated as OSM-derived historical movement signals, not clean per-user histories.

\subsection{Approximate Map Matching}
For each useful pseudo-segment, the system builds a local OSM walking graph, maps GPS points to nearest graph nodes, deduplicates consecutive repeated nodes, and connects successive matched nodes with shortest paths. This approximate map matching is intentionally simple and therefore reported with diagnostics: matched-point coverage, pair success rate, GPS-to-route distance, and route-distance ratio.

For Porto taxi trajectories, the evaluator uses an OSM driving graph and simplifies dense GPS traces into spatially separated anchors before shortest-path stitching. The anchor spacing parameter reduces redundant point-to-point stitching when consecutive samples fall on nearby road nodes. In the current benchmark, an 80 meter minimum anchor spacing and 60-anchor cap modestly reduce route-distance-ratio inflation. This is a reconstruction improvement, but not a full probabilistic map matcher.

\subsection{Route Feature Extraction}
For both reconstructed pseudo-history routes and generated candidate routes, the system extracts OSM-derived features: route distance, major-road percentage, walking-way percentage, residential percentage, service-road percentage, intersections, turns, park proximity, lighting percentage, traffic signals, crossings, tunnel length, and an OSM-only safety proxy. The safety proxy is not a crime model; it is a route-attribute score based on available OSM tags.

\subsection{Dynamic Profile Construction}
Given a request timestamp, the profile module computes temporal context: weekday/weekend, time bucket, season, and rush-hour status. It selects history records with similar context, averages route features, and converts those averages into feature weights. Context rules then adjust weights for rush hour, evening/night, weekend, summer, and winter.

For the Porto benchmark, the generic walking profile is supplemented with a vehicle-aware profile. This profile uses earlier Porto trips as contextual history and emphasizes driving-relevant features: distance, major-road exposure, residential/service-road share, turns, intersections, traffic signals, and tunnel length. The method remains unsupervised: it does not train on held-out labels, and it does not use the oracle candidate during scoring.

\subsection{Candidate Generation}
The route generator constructs a local OSM graph around the query and returns a small set of route alternatives. The walking pipeline uses OSM walking networks and route attributes such as footways, crossings, lighting, and park proximity. The Porto pipeline uses an OSM driving graph and generates diverse vehicle candidates by varying edge weights, including length-oriented, major-road-preferring, major-road-avoiding, turn-simple, and local-road-avoiding variants. Keeping the candidate pool small preserves runtime for the GUI/API workflow, but it also creates an upper-bound limitation: no ranker can recover a path that candidate generation never proposes.

\subsection{Ranking Modes}
The prototype supports five evaluation methods:

\begin{itemize}
  \item \textbf{Random:} random permutation of generated route candidates.
  \item \textbf{Shortest-distance:} candidates sorted by OSM route length.
  \item \textbf{Profile:} candidates scored using dynamic profile weights.
  \item \textbf{Prompt/SBERT:} route summaries ranked by semantic similarity to preference text.
  \item \textbf{Hybrid:} $0.75$ profile score plus $0.25$ prompt/SBERT score, matching the backend.
\end{itemize}

Prompt/SBERT and hybrid modes are skipped when no preference text exists. This is the current case for the OSM pseudo-history records.

\subsection{Scoring}
Profile scoring compares each candidate feature vector with a contextual target profile. For numeric feature $j$, the score rewards closeness between candidate value $x_{ij}$ and target value $\mu_j$ under profile weight $w_j$:

\begin{equation}
s_{\mathrm{profile}}(r_i) = - \sum_j w_j \cdot |\hat{x}_{ij} - \hat{\mu}_{j}|.
\end{equation}

The prompt/SBERT mode converts each route to a textual summary and scores semantic similarity against preference text. The hybrid mode combines normalized profile and prompt scores:

\begin{equation}
s_{\mathrm{hybrid}}(r_i) = 0.75\,s_{\mathrm{profile}}(r_i) + 0.25\,s_{\mathrm{prompt}}(r_i).
\end{equation}

These weights match the backend implementation and are treated as fixed prototype settings, not tuned hyperparameters.
